\chapter{Introduction}\label{ch:intro}

\section{Topic and motivation}

Generative diffusion models have become the dominant paradigm for
synthesising high-dimensional data: they power state-of-the-art systems
for images, audio, video, and molecular structures \citep{sohl2015deep,
ho2020denoising, song2021scorebased, karras2022elucidating}. The same
construction underlies the image and video generators deployed at scale by the
large technology companies, from latent diffusion \citep{rombach2022latent}
onwards, and reaches beyond media: structure prediction in AlphaFold~3 places
atomic coordinates with a diffusion module \citep{abramson2024alphafold3}. Their
central object is the \emph{score function}, the gradient of the
log-density of progressively noised data,
\begin{equation}
  \score(x, t) \;=\; \nabla_x \log p_t(x),
  \label{eq:intro-score}
\end{equation}
which is all a model needs to know in order to reverse a noising process
and turn pure noise back into data \citep{anderson1982reverse}. In
practice the score is approximated by a neural network trained by
denoising score matching \citep{hyvarinen2005estimation,
vincent2011connection}; the network is a black box, and the structure of
the function it has to learn remains poorly understood.

A complementary research programme, initiated in the statistical-physics
community, replaces the black box with \emph{exactly solvable models}:
data distributions simple enough that the score can be computed in
closed form, yet rich enough to exhibit the phenomena observed in real
systems, such as speciation, memorisation, and dynamical phase
transitions along the generation trajectory \citep{biroli2023generative,
biroli2024dynamical, raya2023spontaneous, ghio2024sampling,
achilli2026speciation}. Exact scores act as a microscope: every regime,
transition time, and failure mode can be traced back to explicit
formulas.

This thesis works inside that programme, in a direction that is natural
for data with \emph{temporal or sequential structure}: the data are not
a cloud of independent points but a \emph{trajectory}
$a_0, a_1, \dots, a_{K-1}$ of a Markov process, the prototype of a
video, a time series, or any dynamic object. Each frame of the
trajectory is corrupted independently by an Ornstein--Uhlenbeck (OU)
channel, and the object of study is the \emph{joint} score of the whole
noisy sequence. Two classical bodies of theory then meet the modern one.
First, the joint score of a Gaussian Markov chain is governed by the
algebra of structured matrices: Toeplitz covariances, tridiagonal
precisions, and their common eigenbasis. Second, by an identity that
holds for any data distribution, computing a score is solving a Bayesian
inference problem, and for chain-structured data that problem lives on a
tree-shaped factor graph, the domain of \emph{belief propagation} (BP)
and of the Kalman filter, together with the Gaussian message closures
of statistical physics (TAP, AMP)
\citep{pearl1988probabilistic, mezard2009information, kalman1960new,
donoho2009message}. The thesis develops both routes in full, proves
where they coincide, and measures what is lost when they are forced
apart.

\section{Positioning and gap}

The literature this work builds on can be organised in three strands;
Chapters~\ref{ch:statmech}--\ref{ch:graphs} review them in depth.

\paragraph{Statistical physics of learning machines.} Machine learning
and statistical mechanics share a genealogy that runs from spin glasses
through the Hopfield network and the Boltzmann machine to today's
generative models, a lineage recognised by the 2024 Nobel Prize in
Physics \citep{hopfield1982neural, ackley1985learning, nobel2024physics}.
Its modern branch produced closed-form analyses of diffusion generation
for Gaussian mixtures, ferromagnets, and high-dimensional asymptotics:
the regimes of the backward dynamics and their transitions
\citep{biroli2023generative, biroli2024dynamical}, the general theory of
speciation \citep{achilli2026speciation, sclocchi2025phase}, the
dynamics of memorisation and generalisation in training
\citep{bonnaire2025memorize, garnierbrun2026biased}, and, on the
training side, the demonstration that energy-based models learn through
a cascade of phase transitions \citep{bachtis2024cascade}. These works
treat the data as unstructured collections of samples; the time axis
they study is the diffusion (or training) time. The internal time of a
\emph{sequence}, the fact that a video frame is correlated with its
neighbours through a dynamical rule, is absent from these analyses.

\paragraph{Inference on chains.} Exact posterior inference on
Gauss--Markov chains is classical: the Kalman filter and smoother
\citep{kalman1960new, rauch1965maximum}, hidden Markov models
\citep{rabiner1989tutorial}, and their unifying formulation as
sum--product message passing on factor graphs
\citep{kschischang2001factor, mezard2009information}. What this
literature does not ask is the question that matters for generative
modelling: how does the exact posterior, and therefore the exact score,
deform as a function of the diffusion time, and what does that
deformation imply for anything, algorithm or network, that must compute
it at every noise level?

\paragraph{Message passing with continuous variables.} BP is exact on
trees for discrete variables, where messages are finite probability
vectors. For continuous variables the messages become probability
densities, and practical algorithms must close them within a parametric
family, almost always the Gaussian one; this is the step that connects
BP to AMP and to the TAP equations of statistical physics
\citep{thouless1977solution, donoho2009message, zdeborova2016statistical,
genovese2026amp}. The classical justification of the Gaussian closure is
a central-limit argument valid when every node receives many messages.
On a chain each node receives exactly two: the central-limit argument
fails by construction, and we are not aware of a quantitative assessment
of the Gaussian message closure on chain-structured diffusion
posteriors.

The gap, in one sentence: there is no end-to-end, exactly solvable
account of the joint diffusion score of a structured sequence, covering
its matrix structure, its message-passing computation, and the price of
the Gaussian closure when the sequence is not Gaussian.

\section{Research questions}\label{sec:rq}

The thesis addresses four questions.

\begin{itemize}
  \item[\textbf{RQ1.}] For a stationary Gaussian AR(1) chain observed
  through coordinatewise Ornstein--Uhlenbeck noise, what is the joint
  score, and how do its spatial couplings vary with the diffusion time?

  \item[\textbf{RQ2.}] How can the same Gaussian-chain score be computed
  by belief propagation, and how does the resulting recursion relate to
  Kalman filtering and smoothing?

  \item[\textbf{RQ3.}] For selected non-Gaussian innovation models, how
  accurate is a Gaussian-message approximation relative to a validated
  numerical belief-propagation reference, over the parameter ranges
  studied?

  \item[\textbf{RQ4.}] In the Gaussian model, how rapidly does the
  influence of remote observations on the score decay, and what limited
  guidance does this provide for local approximations of the score?
\end{itemize}

These questions are deliberately confined to the models named in them.
No claim is made about general Markov processes, about universal
properties of learned diffusion scores, or about practical neural
architectures.

\section{Contributions}

\begin{enumerate}
  \item For the stationary Gaussian AR(1) chain under OU corruption, we
  derive the joint score in closed form, diagonalise it for all
  diffusion times in a single eigenbasis, and characterise the evolution
  of its coupling structure: a band-fill law
  $(\Qt)_{i,i+d} = (-1)^{d-1}(2t)^{d-1}(\Qz^d)_{i,i+d} + O(t^d)$ for
  small $t$, a return to isotropy at rate $e^{-2t}$, and exact bulk
  formulas for the posterior variance and the geometric decay rate
  $q(\alpha, t)$ of posterior correlations. The band-fill law corrects a
  prefactor error in an earlier working note of this project; the
  correction is verified numerically (Chapter~\ref{ch:gaussian}).

  \item We derive the belief-propagation computation of the same score
  on the chain's factor graph, show algebraically that Gaussian messages
  are closed under the updates, identify the resulting recursion, update
  by update, with the Kalman filter and the Rauch--Tung--Striebel
  smoother, and verify that it reproduces the matrix score to within
  double-precision rounding across a parameter grid
  (Chapter~\ref{ch:gaussian}).

  \item For Laplace-innovation chains we establish which identities
  remain exact, locate the precise step at which closed forms stop, and
  measure the accuracy of Gaussian-projected BP against a numerically
  validated grid implementation of the functional recursion: the median
  relative score error is $0.39$ at $t = 0.02$ and decays below
  $10^{-3}$ for $t \gtrsim 0.9$, with the score direction substantially
  more accurate than its magnitude, in the tested regimes
  (Chapter~\ref{ch:laplace}).

  \item In the Gaussian model we prove that the error of a radius-$r$
  local estimator of the score decays as $q(\alpha,t)^r$, and we measure
  that truncating the inference (windowed message passing) is more
  accurate than truncating the score matrix (banding) at small diffusion
  times (Chapter~\ref{ch:gaussian}).

  \item As a separate, self-contained result, we characterise the
  per-node (AMP/TAP) closure on the same model: it returns the exact
  score but a different posterior variance, with an existence boundary
  in closed form, a breakdown time $t_c(\alpha)$, and a critical
  coupling $\alpha_c = \sqrt{2} - 1$. All formulas are independently
  re-verified; the material is placed in Appendix~\ref{app:amp} because
  it branches off the main line (statement of results in
  Section~\ref{sec:conclusions-amp}).
\end{enumerate}

\section{Methodology}

The methodology is deliberate: no neural networks, no black boxes;
closed forms first, numerics as audit. Every claim is either proved by
explicit algebra (with the longer passages isolated in numbered
\emph{toolboxes}) or verified by deterministic numerical checks against
ground truth computed by an independent method, in the spirit of the
exactly solvable programme of \citet{biroli2023generative}. Where an
approximation is introduced (the Gaussian closure of BP messages, local
truncations), its error is \emph{measured} against an exact reference
(matrix algebra, or dense grid quadrature validated on the solvable
case). All numbers in the thesis are regenerated by deterministic audit
scripts, documented in Appendix~\ref{app:repro}; we regard this
reproducibility as a strength of the method, not as a contribution in
itself.

\section{Scope and limitations}

The exact results concern one-dimensional state per frame, linear AR(1)
dynamics, and stationary initial conditions. The non-Gaussian results
are numerical, valid within the validated resolution regime of the
reference solver and the tested parameter ranges. The locality results
are proved for the Gaussian bulk only. No learning experiments are
performed, so any implication for trained score networks is a hypothesis
suggested by a controlled model, not an established property of
practical systems.

\section{Structure of the thesis}

Chapters~\ref{ch:statmech}--\ref{ch:graphs} contain the background and
literature review; Chapters~\ref{ch:development}--\ref{ch:em-results}
contain the research. Throughout, each concept is developed once, in the
chapter where it naturally belongs, and then used without repetition.

Chapter~\ref{ch:statmech} develops the statistical mechanics this thesis
uses (phase space and Liouville's theorem, the Boltzmann--Gibbs
distribution, free energy, the Ising model) and then follows the
discipline into machine learning: spin glasses, the Hopfield network,
the Boltzmann machine and its restricted variant, and the statistical
mechanics of learning. Chapter~\ref{ch:stochproc} supplies the
probabilistic dynamics: stochastic processes and stationarity, Markov
chains, the AR(1) running example, sampling by dynamics, the It\^o
calculus, the Ornstein--Uhlenbeck corruption channel, the Fokker--Planck
equation, and Anderson's time-reversal theorem.
Chapter~\ref{ch:diffusion} reviews modern generative modelling:
energy-based models in LeCun's formulation, score matching, the DDPM
construction, the score-SDE synthesis, flow models, the score--posterior
identity, the statistical-physics analyses of diffusion, and the cascade
of phase transitions in energy-based-model training.
Chapter~\ref{ch:graphs} builds the inference machinery: the multivariate
Gaussian and its precision-matrix reading, the spectral theorem, Markov
random fields, factor graphs, the sum--product algorithm, Gaussian BP
and the Kalman filter, the TAP/AMP closures, and the architectural
mirror of locality in convolutional networks and U-Nets.

Chapter~\ref{ch:development} records the preliminary work: the toy
models, the first supervision models, and the corrections that produced
the final research questions. Chapter~\ref{ch:gaussian} solves the
Gaussian chain exactly (RQ1, RQ2, and the Gaussian half of RQ4), with
the spectral form of the precision matrix as the organising object and
belief propagation identified with the Kalman smoother.
Chapter~\ref{ch:laplace} leaves Gaussianity in the most controlled way
possible, Laplace innovations, and measures what the Gaussian message
family costs when the data are no longer Gaussian (RQ3).
Chapter~\ref{ch:ring} extends the state at each site from a scalar to a
point on a rotating ring, a nonlinear dynamic object, and shows that a
Cartesian surrogate is exactly solvable while the model as drawn is
solvable numerically to a measured accuracy, both admitting the same
structural description of the joint score.
Chapters~\ref{ch:em-method}--\ref{ch:em-results} leave the innovation law
itself unknown: an EM procedure built on the exact chain inference of
Chapter~\ref{ch:gaussian} estimates it from noised data alone, and the
structured estimator is compared against trained neural denoisers, both
pointwise and generatively, along axes of model capacity and of
misspecification of the Markov assumption itself.
Chapter~\ref{ch:conclusions} states findings per research question,
limitations, and the questions that follow. Three appendices collect the
Gaussian identities with proofs, the AMP/TAP analysis with its exact
phase boundary, and the reproducibility protocol.
